\documentclass{article}

\usepackage[accepted]{icml2024}
\usepackage{microtype}
\usepackage{graphicx}
\usepackage{booktabs}
\usepackage{amsmath}
\usepackage{amssymb}
\usepackage{hyperref}

\newcommand{\theHalgorithm}{\arabic{algorithm}}
\icmltitlerunning{Layer-wise LoRA Ranks and Catastrophic Forgetting}

\InputIfFileExists{generated/results.tex}{}{
  \PackageError{rank-lora-paper}
  {Missing generated/results.tex}
  {Run the complete experiment matrix and then python scripts/build_paper.py.}
}

\begin{document}

\twocolumn[
\icmltitle{Do Transformer Layers Prefer Different LoRA Ranks?\\
A Controlled Study of Sequential Fine-Tuning and Forgetting}

\begin{icmlauthorlist}
\icmlauthor{Matheus Costa Monteiro}{independent}
\end{icmlauthorlist}
\icmlaffiliation{independent}{Independent Researcher}
\icmlcorrespondingauthor{Matheus Costa Monteiro}{matheuscostamonteiro.mc@gmail.com}
\vskip 0.3in
]

\begin{abstract}
Low-rank adaptation usually assigns the same rank to every adapted Transformer
layer, although layer depth may induce different capacity requirements. We
study whether reallocating a fixed LoRA parameter budget across depth changes
both task acquisition and catastrophic forgetting. QLoRA adapters on
Qwen3-4B-Instruct are trained sequentially on instruction following and
mathematical reasoning in both task orders. Uniform, early-heavy, middle-heavy,
late-heavy, and data-driven spectral allocations are compared at equal
parameter count. \GeneratedAbstractResult
\end{abstract}

\section{Introduction}
LoRA freezes pretrained weights and represents each adapted linear update as
$\Delta W_\ell=B_\ell A_\ell$, where rank $r_\ell$ controls trainable capacity
\citep{hu2022lora}. A uniform $r_\ell$ is convenient but need not be optimal:
early, intermediate, and late Transformer blocks play different computational
roles. Adaptive methods such as AdaLoRA and dEBORA allocate capacity across
modules \citep{zhang2023adalora,zangrando2025debora}, but their interaction
with sequential learning remains underexplored.

We isolate allocation from parameter count and scaling. Every experiment adapts
only \texttt{q\_proj} and \texttt{v\_proj}, uses conventional LoRA with
$\alpha_\ell=r_\ell$, and matches the rank-16 budget within one percent. We ask
(i) whether the location of rank affects acquisition and retention, (ii)
whether the effect depends on task order, and (iii) whether a low-cost spectral
pilot improves on uniform ranks.

\section{Experimental protocol}
\paragraph{Tasks and contamination control.}
Instruction following is trained on the filtered IFEval-like corpus and
evaluated on the official IFEval benchmark using prompt-level strict accuracy
\citep{zhou2023ifeval}. Mathematical reasoning is trained on NuminaMath-CoT
\citep{li2024numinamath} and evaluated by numerical exact match on GSM8K
\citep{cobbe2021gsm8k}. Before splitting, benchmark-source rows, normalized
exact hashes, and approximate character-$5$-gram matches above 0.90 are
removed. Training and calibration are disjoint and contain no evaluation
prompts. Each task contributes exactly one million supervised assistant tokens.

\paragraph{Sequential QLoRA.}
We use NF4 4-bit QLoRA \citep{dettmers2023qlora}, BF16 computation, gradient
checkpointing, a learning rate of $2{\times}10^{-4}$, and an effective batch
size of 16. The same adapter is trained on task 1, saved, reloaded, and
continued on task 2; optimizer and scheduler are reset between tasks. Both
orders are evaluated at the base, stage-1, and stage-2 checkpoints.

\paragraph{Rank allocations.}
The 36 blocks are partitioned into 12 early, 12 middle, and 12 late blocks.
Uniform ranks are $(16,16,16)$; the three manual allocations are permutations
of $(32,8,8)$. All contain \NumLoraParameters trainable LoRA parameters.
The spectral pilot trains rank 32 on an equal-token task mixture. Singular
values of $BA$ are computed from QR factorizations of $B$ and $A^\top$ and an
SVD of the resulting $32{\times}32$ matrix. Marginal spectral energy then
allocates ranks from $\{4,8,\ldots,32\}$ under the same budget.

\section{Results}
\begin{figure}[t]
  \centering
  \includegraphics[width=\columnwidth]{../results/figures/performance_by_stage.pdf}
  \caption{Benchmark performance at the base, stage-1, and stage-2 checkpoints.
  All values are on a 0--100 scale.}
  \label{fig:stages}
\end{figure}

\begin{figure}[t]
  \centering
  \includegraphics[width=\columnwidth]{../results/figures/forgetting_by_configuration.pdf}
  \caption{Forgetting by rank allocation and benchmark. Lower is better.}
  \label{fig:forgetting}
\end{figure}

We define forgetting as the best observed score minus the final score, backward
transfer as final minus immediate post-learning performance, and net gain as
final minus base. Figure~\ref{fig:stages} exposes acquisition and retention
without mixing IFEval and GSM8K into a single headline metric.
\GeneratedMainResult

The manual allocation promoted to three seeds was selected only through
assistant-token negative log-likelihood on decontaminated calibration data.
The selected configuration was \SelectedManualConfig; neither IFEval nor GSM8K
participated in that decision. \GeneratedOrderResult

\section{Conclusion}
\GeneratedConclusion
These findings characterize a budget-allocation effect rather than proving
that non-uniform ranks are universally superior. Future work should test more
task families, model scales, and longer task sequences.

\clearpage
\bibliography{references}
\bibliographystyle{icml2024}

\appendix
\section{Reproducibility details}
The released pipeline pins model and dataset revisions, stores the complete
rank pattern, logs parameters and peak memory per GPU, saves all generations,
and supports resumption at each training stage. Scores are generated greedily
with temperature zero. The initial sweep uses seed 42; uniform, the selected
manual allocation, and the spectral allocation use seeds 7, 42, and 123.

\section{Additional diagnostics}
\begin{figure}[h]
  \centering
  \includegraphics[width=.8\textwidth]{../results/figures/rank_by_depth.pdf}
  \caption{Applied rank as a function of Transformer depth.}
\end{figure}
\begin{figure}[h]
  \centering
  \includegraphics[width=.8\textwidth]{../results/figures/spectral_energy_curves.pdf}
  \caption{Cumulative spectral-energy curves from the discarded pilot adapter.}
\end{figure}

\end{document}

